Use the actual odd form and unaveraged point of
Propositions~\ref{ex:app-theta-form} and~\ref{ex:app-theta-comparison}.
There is one constant $C_{\rm det}$, independent of the actual stage
and admissible $D$, for which (O.8) holds. Explicitly, for a known
simple-zero $D$,
$v_2a_M(D)\ge\min\{M,j_{\rm gen}(D)-C_{\rm det}\}$.
For every admissible $D$, the strict inequality $v_2a_M(D)<M$
implies $\operatorname{ord}_{s=1}L(E^D,s)=1$ and
$j_{\rm gen}(D)\le v_2a_M(D)+C_{\rm det}$.
If an admissible odd-sign $D$ is not a simple zero, its coefficient
vanishes modulo $2^M$. No genus index of a torsion point is defined
or used in this assertion.
\end{proposition}
\begin{proof}\leavevmode
For a known simple-zero $D$, write $P_D=m_DG+T$ in its native
rank-one free lattice. The same $L$ kills $T$.
An integral finite-field Kummer map cannot decrease
$v_2(m_D)$, so (O.7) gives the lower bound by the native index.
The rational free line has index at most two in that native
lattice, by conjugation and norm. Absorb this and the fixed
initial convention in $C_{\rm det}$. The first inequality below
holds for a known simple-zero $D$; the visibility implication holds
for every admissible $D$:
\[
 \begin{split}
 v_2a_M(D)&\ge\min\{M,j_{\rm gen}(D)-C_{\rm det}\},\\
 v_2a_M(D)<M&\ \Longrightarrow\
  \ord_{s=1}L(E^D,s)=1,\quad
       j_{\rm gen}(D)\le v_2a_M(D)+C_{\rm det}.
 \end{split}                                                    \tag{O.8}
\]
If an odd-sign $D$ is not a simple zero, the height formula gives a
torsion genus point, killed by $L$, and its coefficient is zero
modulo $2^M$. Combined with the first inequality on simple-zero
parameters, this gives the all-$D$ truncated bound used below,
without assigning a rank-one lattice index to a torsion point.
\end{proof}

